\documentclass[11pt]{article}
\usepackage{mathblatt}
% Präambel des Setzers (werkzeuge/setzer-praeambel.tex, wortgleich); dazu \gteil (Muster T2U)
\makeatletter
\setlength{\mbpfbe}{0mm}% keine Punkte (pruefung.md § 5: nur beim Prüfstein)
\newlength{\gEin}\setlength{\gEin}{0pt}
\let\gSchrift\relax
\renewcommand{\pfkreuzzeile}[1]{\par\noindent\hangindent3.2em\hangafter1\hspace*{1.6em}$\square$\ #1\par}
\newcommand{\gkreuz}[1]{$\square$\ #1\hspace{2em}}
\newcommand{\gkreuzzeile}[2]{\par\noindent\hangindent2.2em\hangafter1\mbox{#1\,$\square$}\ #2\par}% Buchstabe, Option: eine je Zeile (neben Skizze)
% links, was man ansieht, rechts, was man tut; Spaltengrenze überall an derselben Stelle (gemeinsam.md „Satz“)
% v0.17: Fußeintrag „Nr: Wert“ – kurz in einer Box (nie getrennt), lang umbrechbar mit der Nummer am Anfang
% (vorher \mbox: lange Lösungen liefen über den Rand, AS40 138 pt); robust, weil \fusshilfe in die .aux schreibt
\NewDocumentCommand{\gfuss}{m m}{\sbox0{#1:~#2}\ifdim\wd0<0.45\textwidth\mbox{#1:~#2}\else\mbox{#1:}\nobreak\ #2\fi}
\newsavebox{\gbox}\newlength{\gLinks}
% v0.10 (Lehrer 13.10.): keine starre Halbierung – die linke Spalte ist so breit wie die Figur (+ 6 mm), die
% Antwort steht direkt daneben; ist die Figur breiter als 0,62 Zeile, steht die Antwort darunter
\newcommand{\gzwei}[2]{\par\smallskip\noindent\sbox{\gbox}{#1}%
  \setlength{\gLinks}{\dimexpr\wd\gbox+6mm\relax}%
  \ifdim\gLinks>0.62\linewidth
    \usebox{\gbox}\par\smallskip\noindent\begin{minipage}{\linewidth}#2\end{minipage}%
  \else
    \begin{minipage}[t]{\gLinks}\vspace{0pt}\usebox{\gbox}\end{minipage}%
    \begin{minipage}[t]{\dimexpr\linewidth-\gLinks\relax}\vspace{0pt}\raggedright #2\end{minipage}%
  \fi\par}
% v0.9 (12.10., pruefung.md § 5 „Marke mit B, Bedingungen rechts“): drittes Argument von pfaufg = Bedingungen
% der Prüfung (v0.10: \pfohnetr = TR durchgestrichen, „FOR“), klein grau im rechten Rand auf der Grundlinie der Nummer, Zeilen mit \\
\newcommand{\gbedrand}[2]{\ifblank{#2}{}{\rlap{\hspace*{\dimexpr#1+2mm\relax}%
  \parbox[t]{15mm}{\raggedright\scriptsize\leavevmode\color{mbgrau}\strut#2}}}}
% Teil eines Bündels (a, b …): Marke links im Rand, Bedingungen rechts im Rand, auf der Zeile des Teils
\newcommand{\gteilrand}[2]{\llap{\mbpfrandmarke{#1}\hspace*{\mbpfnr}}\gbedrand{\linewidth}{#2}}
% Teil eines Bündels: Buchstabe in eigener Spalte, Inhalt eingerückt (\gEin = Einzug, damit die Spaltengrenze von
% \gzwei an derselben Stelle bleibt)
\newenvironment{gbteil}[3]{\par\smallskip\noindent\gteilrand{#2}{#3}\makebox[\mbpfteil][l]{\textbf{#1}}%
  \begin{minipage}[t]{\dimexpr\linewidth-\mbpfteil\relax}\setlength{\gEin}{\mbpfteil}\raggedright\strut\ignorespaces}%
 {\end{minipage}\par\addvspace{6pt}}% v0.18 (PG80 Nr. 1): Abstand, damit die Linie von a) nicht an b) rutscht
% Grafik höchstens so breit wie die Spalte (Bündel nebeneinander, Koordinatensystem ohne Karo)
\newsavebox{\gpbox}
\newcommand{\gpassend}[1]{\sbox{\gpbox}{#1}\ifdim\wd\gpbox>\linewidth\typeout{GPASSEND: verkleinert \the\wd\gpbox\space auf \the\linewidth}\resizebox{\linewidth}{!}{\usebox{\gpbox}}\else\usebox{\gpbox}\fi}
% Skizzenfeld (v0.10, Lehrer 13.10.): Karo ab dem linken Rand über die Zeile, „Skizze“ klein grau oben links im Feld
% (wie \pfskizzenfeld der Vorlage); der rechte Rand bleibt den Bedingungen
\newcommand{\gskizze}[1]{\par\smallskip\noindent\begin{tikzpicture}
  \clip (0,0) rectangle ({\linewidth-1mm},{#1*10mm});
  \draw[black!16,line width=0.3pt,step=5mm] (0,0) grid ({\linewidth},{#1*10mm});
  \node[anchor=north west,inner sep=2pt,font=\scriptsize,text=mbgrau,fill=white] at (0,{#1*10mm}) {Skizze};
  \end{tikzpicture}\par}
\RenewDocumentEnvironment{pfaufg}{m m m}{%
  \begin{lrbox}{\mbaufgabebox}\begin{minipage}[t]{\linewidth}%
  \gSchrift\noindent\hspace*{\gEin}\mbpfrandmarke{#1}\gbedrand{\linewidth-\gEin}{#3}%
  \begin{minipage}[t]{\mbpfnr}\strut\textbf{#2}\end{minipage}%
  \begin{minipage}[t]{\dimexpr\linewidth-\gEin-\mbpfnr\relax}\raggedright\strut\ignorespaces}%
 {\end{minipage}%
  \end{minipage}\end{lrbox}%
  \setlength{\mbaufgabehoehe}{\dimexpr\ht\mbaufgabebox+\dp\mbaufgabebox\relax}%
  \ifdefempty{\gkopfnach}{\par\addvspace{5pt}\Needspace*{\mbaufgabehoehe}}{%
    % Stufenkopf, der mit dieser (ersten) Aufgabe zusammenbleibt: Platz für beide verlangen, dann den Kopf setzen
    % (12.10., Quadratische: ein Kopf stand allein unten, die Aufgabe mit Koordinatensystem auf der nächsten Seite)
    % v0.17 (AS40 „Terme deuten“ allein unten): hinter dem Kopf steht das Sprungziel (\special, ein Knoten, nach
    % dem die Seite brechen durfte) – darum \nobreak nach dem Kopf; Vorhalt 56 pt = Kopf mit Linie und Abständen
    \par\Needspace*{\dimexpr\mbaufgabehoehe+56pt\relax}\gkopfnach\par\nobreak\gdef\gkopfnach{}\addvspace{5pt}}%
  \noindent\usebox{\mbaufgabebox}\par\addvspace{5pt}}
% v0.18 (plan.md § 10, 15.10. l): Abschnittsbuchstabe A, B, C vor dem Namen im Kopf – als Vorsatz, damit der
% Name im TeX gleich bleibt (Bestellkarte, Sprungziele lesen ihn)
\def\gbstakt{}\newcommand{\gbstvor}[1]{\gdef\gbstakt{#1}}
\NewCommandCopy\gpfstufekopfalt\pfstufekopf
\RenewDocumentCommand{\pfstufekopf}{o m m}{\IfValueTF{#1}{\gpfstufekopfalt[#1]}{\gpfstufekopfalt}%
  {\ifx\gbstakt\empty\else\gbstakt\enspace\fi #2}{#3}\gdef\gbstakt{}}
% v0.18 (16.10. e): Begriff grau, ohne Nummer, ohne Rahmen über der ersten Aufgabe der Stufe
\newcommand{\gbegriff}[1]{\par\nobreak\noindent\hspace*{\mbpfnr}\parbox[t]{\dimexpr\linewidth-\mbpfnr\relax}{\raggedright\small\color{mbgrau}#1}\par\nobreak\addvspace{3pt}}
% v0.18 (16.10. g): Lösungsblatt – Kopf mit kräftigem Strich, Abschnitt klein grau als Zwischenzeile
\newcommand{\glkopf}[1]{\noindent{\large\bfseries Lösungen \textperiodcentered{} #1}\par\vspace{3pt}\hrule height 1.2pt\par\medskip}
\newcommand{\glzwischen}[1]{\par\addvspace{8pt}\Needspace*{4\baselineskip}\noindent{\footnotesize\color{mbgrau}#1}\par\nobreak\vspace{1pt}}
% v0.18 (15.10. o, w): Blattende – Kern ohne Rahmen, dann graue Zeile „Mehr: n Aufgaben“
\newcommand{\gkern}[1]{\par\addvspace{14pt}\Needspace*{6\baselineskip}\noindent{\bfseries Kern}\par\nobreak\smallskip
  {\small\raggedright #1\par}}
\newcommand{\gmehr}[1]{\par\addvspace{10pt}\noindent{\footnotesize\color{mbgrau}Mehr: #1 Aufgaben}\par}
% Stufenkopf vormerken: \gkopfvor{<Kopf>} – die nächste pfaufg setzt ihn, sobald beide zusammen Platz haben
\def\gkopfnach{}
\newcommand{\gkopfvor}[1]{\gdef\gkopfnach{#1}}
% Teilaufgabe mit Skizze links in gzwei: Buchstabe in eigener Spalte, darunter Text und Skizze
\newcommand{\gteil}[2]{\makebox[\mbpfteil][l]{\textbf{#1}}\parbox[t]{\dimexpr0.5\textwidth-\gEin-\mbpfnr-5mm-\mbpfteil\relax}{\raggedright #2}}
% Stufenkopf der Vorlage (mathblatt.sty 2026-10-12: Linie, Abstand, 13 pt, \pfjahre) ohne ihren Vorhalt von
% 14 Zeilen: der Setzer fordert vor dem Kopf 0,16 Seitenhöhe (Kopf + erste Aufgabe), mehr nicht
\setlength{\mbpfstufevorhalt}{0pt}
\makeatother
\newcommand{\karo}[1]{\par\smallskip\noindent\begin{tikzpicture}
  \clip (0,0) rectangle ({\linewidth-1mm},{#1*10mm});
  \draw[black!16,line width=0.3pt,step=5mm] (0,0) grid ({\linewidth},{#1*10mm});
  \end{tikzpicture}\par}
\newcommand{\kk}[1]{$\square$\,#1\hspace{0.9em}}
\newcommand{\linie}[1]{\,{\color{black!45}\rule[-1pt]{#1}{0.4pt}}\,}
\newcommand{\pkt}{\textperiodcentered{}}
% v0.14 (14.10.): Stufenkopf mit fünf Punkten (Option = Jahre für die Bestellkarte) oder „zuletzt 19“
\NewDocumentCommand{\gpunkte}{o m}{#2}
\newcommand{\gzuletzt}[1]{zuletzt #1}
% v0.7 (12.10.): Schülerfrage unter dem Stufenkopf; Ankreuzoption als Bild (Buchstabe und Kästchen darunter)
\newcommand{\gfrage}[1]{\par\nobreak\noindent\hspace*{\mbpfnr}{\small\itshape #1}\par\nobreak\addvspace{2pt}}
\newlength{\gbildbreite}\setlength{\gbildbreite}{0.32\linewidth}% Setzer setzt sie je Reihe (vier Bilder: 0,24)
\newcommand{\gbildoption}[3]{\begin{minipage}[t][#3][b]{\gbildbreite}\centering #2\par\smallskip\mbox{#1\,$\square$}\end{minipage}}
% Lineare (12.10.): Bild mit Buchstabe ohne Kästchen (Diagramm zuordnen); Aussage mit „wahr / falsch“ zum Ankreuzen
\newcommand{\gbildname}[3]{\begin{minipage}[t][#3][b]{\gbildbreite}\centering #2\par\smallskip\mbox{\textbf{#1}}\end{minipage}}
\newcommand{\gwahrfalsch}[2]{\par\noindent\hangindent2.2em\hangafter1\makebox[2.2em][l]{(#1)}#2\hfill\mbox{wahr $\square$\quad falsch $\square$}\par}
\begin{document}

\pfheftstil
\pflaufkopf[Abitur GK \pkt{} Prüfung]{Tangente}{AA40}\pfseitevon
\pfheftkopf{Tangente}{}
\pfstufenzeile{Tangentengleichung gesucht \pkt{} Tangente nachweisen \pkt{} Berührpunkt gesucht}

% ===================== Stufe: Tangentengleichung gesucht =====================
\gkopfvor{\pfstufekopf{Tangentengleichung gesucht}{\gpunkte[26 25 23 22]{\textcolor{black!80}{$\bullet$}\,\textcolor{black!80}{$\bullet$}\,\textcolor{black!25}{$\circ$}\,\textcolor{black!80}{$\bullet$}\,\textcolor{black!80}{$\bullet$}}}\special{pdf:dest (stufe-tangentengleichung-gesucht) [@thispage /XYZ @xpos @ypos null]}}
% kein Einstieg: keine Schülerfrage in der Gliederung (Schülerfrage, -Fälle, -Form, -Weg, -Skizze; oder Frage) – Einstieg nur von dort, nie aus der Bank (pruefung.md § 5, 12.10.)
% 2022-bebb-gk-A1.1b-m1 (Original 2022-bebb-gk-A1.1b, Schwere I, rang 11; herausgelöste Fassung: keine)
\begin{pfaufg}{GK~’22\\A1.1b}{1.}{\pfohnetr}
Gegeben ist die Funktion $f$ mit $f(x) = x^3 - 10x^2 + 21x$. Bestimme eine Gleichung der Tangente an den Graphen von $f$ im Punkt $P(1 \mid f(1))$.
\karo{4}
\fusshilfe{\gfuss{1}{4x + 8}}
\end{pfaufg}
% 2025-bebb-gk-A1.1a-m1 (Original 2025-bebb-gk-A1.1a, Schwere I, rang 12; herausgelöste Fassung: keine)
\begin{pfaufg}{GK~’25\\A1.1a}{2.}{\pfohnetr}
Gegeben ist die in $\mathbb{R}$ definierte Funktion $f$ mit $f(x) = \tfrac{1}{3}x^3 - x^2 + 2$; die Abbildung zeigt ihren Graphen $G$ und den Punkt $P(3 \mid 2)$. Bestimme rechnerisch eine Gleichung der Tangente $t$ an $G$ in $P$.
\gzwei{\begin{abiksys}[xmin=-3,xmax=5,ymin=-5,ymax=6,le=0.5]\funktion[4.1]{\x^3/3-\x^2+2}{G}\punkt{3}{2}{P}\end{abiksys}}{\karo{4}}
\fusshilfe{\gfuss{2}{3x $-$ 7}}
\end{pfaufg}
% 2026-bb-gk-B2.2b-m1 (Original 2026-bb-gk-B2.2b, Schwere I, rang 13; herausgelöste Fassung: keine)
\begin{pfaufg}{GK~’26\\B2.2b}{3.}{}
Gegeben ist die Funktion $f$ mit $f(x) = 2e^x + 3$, $x \in \mathbb{R}$; ihr Graph ist $G_f$. Ermittle eine Gleichung der Tangente $t$ an $G_f$ im Schnittpunkt von $G_f$ mit der $y$-Achse (zur Kontrolle: $y = 2x + 5$).
\karo{4}
\fusshilfe{\gfuss{3}{2x + 5}}
\end{pfaufg}
% 2023-bebb-gk-A1.1a-m1 (Original 2023-bebb-gk-A1.1a, Schwere II, rang 14; herausgelöste Fassung: keine)
\begin{pfaufg}{GK~’23\\A1.1a}{4.}{\pfohnetr}
Der Graph einer Funktion $f$ ist symmetrisch zur $y$-Achse. Die Tangente $t_1$ an den Graphen im Punkt $(1 \mid f(1))$ hat die Gleichung $y = \tfrac{3}{4}x + 3$. Gib eine Gleichung der Tangente $t_2$ im Punkt $(-1 \mid f(-1))$ an und begründe deine Angabe.
\karo{4}
\fusshilfe{\gfuss{4}{$-\tfrac{3}{4}$ x + 3}}
\end{pfaufg}

% ===================== Stufe: Tangente nachweisen =====================
\gkopfvor{\pfstufekopf{Tangente nachweisen}{\gpunkte[26 23]{\textcolor{black!80}{$\bullet$}\,\textcolor{black!25}{$\circ$}\,\textcolor{black!25}{$\circ$}\,\textcolor{black!80}{$\bullet$}\,\textcolor{black!25}{$\circ$}}}\special{pdf:dest (stufe-tangente-nachweisen) [@thispage /XYZ @xpos @ypos null]}}
% kein Einstieg: keine Schülerfrage in der Gliederung (Schülerfrage, -Fälle, -Form, -Weg, -Skizze; oder Frage) – Einstieg nur von dort, nie aus der Bank (pruefung.md § 5, 12.10.)
% 2026-bb-gk-B2.1d-m1 (Original 2026-bb-gk-B2.1d, Schwere I, rang 15; herausgelöste Fassung: keine)
\begin{pfaufg}{GK~’26\\B2.1d}{5.}{}
Gegeben ist die in $\mathbb{R}$ definierte Funktion $f$ mit $f(x) = \tfrac{2}{3}x^4 + \tfrac{4}{3}x^3 + 2$ und ihrer Ableitung $f'(x) = \tfrac{4}{3}x^2 \cdot (2x + 3)$. Weise nach, dass die Gerade $t$ mit $t(x) = \tfrac{4}{3}x + \tfrac{8}{3}$ den Graphen von $f$ im Punkt $\left(-1 \mid \tfrac{4}{3}\right)$ berührt.
\karo{4}
\fusshilfe{\gfuss{5}{$\tfrac{4}{3}$; t ist Tangente in ($-$1|$\tfrac{4}{3}$)}}
\end{pfaufg}
% 2023-bebb-gk-A1.3b-m1 (Original 2023-bebb-gk-A1.3b, Schwere II, rang 16; herausgelöste Fassung: keine)
\begin{pfaufg}{GK~’23\\A1.3b}{6.}{\pfohnetr}
Gegeben sind die Funktionen $f$ mit $f(x) = x^3 - 6x^2$ und $g$ mit $g(x) = 15x + 8$ sowie der Punkt $P(-1 \mid -7)$. Begründe, dass sich die Graphen von $f$ und $g$ im Punkt $P$ berühren.
\karo{4}
\fusshilfe{\gfuss{6}{die Graphen berühren sich}}
\end{pfaufg}

% ===================== Stufe: Berührpunkt gesucht =====================
\gkopfvor{\pfstufekopf{Berührpunkt gesucht}{\gpunkte[25 23]{\textcolor{black!25}{$\circ$}\,\textcolor{black!80}{$\bullet$}\,\textcolor{black!25}{$\circ$}\,\textcolor{black!80}{$\bullet$}\,\textcolor{black!25}{$\circ$}}}\special{pdf:dest (stufe-beruhrpunkt-gesucht) [@thispage /XYZ @xpos @ypos null]}}
% kein Einstieg: keine Schülerfrage in der Gliederung (Schülerfrage, -Fälle, -Form, -Weg, -Skizze; oder Frage) – Einstieg nur von dort, nie aus der Bank (pruefung.md § 5, 12.10.)
% 2025-bebb-gk-A1.1b-m1 (Original 2025-bebb-gk-A1.1b, Schwere II, rang 17; herausgelöste Fassung: keine)
\begin{pfaufg}{GK~’25\\A1.1b}{7.}{\pfohnetr}
Gegeben ist die Funktion $f$ mit $f(x) = \tfrac{1}{3}x^3 - x^2 + 2$; die Abbildung zeigt ihren Graphen $G$. Die Tangente an $G$ im Punkt $P(3 \mid 2)$ hat die Steigung 3. Es gibt genau eine weitere Tangente an $G$, die zu ihr parallel ist. Bestimme ihren Berührpunkt und skizziere diese Tangente in der Abbildung.
\par\smallskip\noindent\begin{abiksys}[xmin=-3,xmax=5,ymin=-5,ymax=6,le=0.5]\funktion[4.1]{\x^3/3-\x^2+2}{G}\punkt{3}{2}{P}\end{abiksys}\par
\karo{2}
\fusshilfe{\gfuss{7}{Berührpunkt ($-$1|$\tfrac{2}{3}$), Tangente mit Steigung 3 durch diesen Punkt}}
\end{pfaufg}
% 2023-bebb-gk-B2.2d-m1 (Original 2023-bebb-gk-B2.2d, Schwere II, rang 18; herausgelöste Fassung: keine)
\begin{pfaufg}{GK~’23\\B2.2d}{8.}{}
Gegeben ist die Funktion $f$ mit $f(x) = x^3 - 9x^2 + 24x - 16$, $x \in \mathbb{R}$. Die Tangente an den Graphen von $f$ im Punkt $S(0 \mid f(0))$ hat mit dem Graphen einen weiteren gemeinsamen Punkt. Bestimme seine Koordinaten.
\karo{4}
\fusshilfe{\gfuss{8}{24x $-$ 16; weiterer Punkt (9|200)}}
\end{pfaufg}

\pfblattende{10}{26 & 25 & 24 & 23 & 22 & 21 & 20 & 19 & 18 & 17}{B2.1d B2.2b & A1.1a A1.1b &  & A1.1a A1.3b B2.2d & A1.1b &  & B2.2h & B2.2f &  & }%
{\pfausschnitt{Änderungsrate}{Tangente}{Winkel am Graphen \pkt{} Dreieck aus Tangente und Normale}}
\end{document}